% !TeX root = ../main.tex

\appendix

\chapter{Supplementary Material}\label{chapter:appendix}

\section{Detailed Bill of Materials}\label{sec:bom-detailed}

This section provides a complete component list for the experimental X-wing quadrotor platform, including part numbers, specifications, and suppliers for reproducibility.

\begin{table}[h]
\centering
\caption{Complete Bill of Materials (BOM) for the baseline experimental platform.}
\label{tab:bom-detailed}
\small
\begin{tabularx}{\textwidth}{ll>{\raggedright\arraybackslash}X}
\toprule
Category & Component & Model / Key Specs \\
\midrule
\multicolumn{3}{l}{\textbf{Airframe}} \\
Frame & Center frame & X-wing center frame, 3D-printed PLA \\
Frame & Wing spars & Carbon fiber tubes, 10\,mm OD, 8\,mm ID \\
Wings & Wing skins & Four fixed wings (length 450\,mm, chord 250\,mm; NACA~0015 airfoil), orthogonal layout ($90^{\circ}$ between adjacent wings); 3D-printed Bambu Lab PLA Aero; total projected horizontal area $S_\mathrm{t}\approx 0.318$\,m$^2$ \\
\midrule
\multicolumn{3}{l}{\textbf{Propulsion}} \\
Motors & T-MOTOR VELOX V2808 ($\times$4) & KV1300, max thrust $\approx$22\,N each @ 6S \\
Propellers & HQProp 7×3.5×3 ($\times$4) & 3-blade racing propeller (7", light grey) \\
\midrule
\multicolumn{3}{l}{\textbf{Avionics}} \\
Flight Controller & Kakute H7 v1.5 & STM32H743, ICM-42688-P IMU @ 8\,kHz \\
ESC & Holybro Tekko32 F4 4in1 & 4-in-1 ESC, 65A, DShot600/1200 capable \\
\midrule
\multicolumn{3}{l}{\textbf{Power}} \\
Battery ($\times$2, parallel) & Tattu R-Line V5.0 6S 1400\,mAh & 150C discharge, XT60 connector; effective 6S 2800\,mAh (parallel) \\
Companion Power & Tattu R-Line V3.0 4S 1800\,mAh & 120C discharge, XT60 connector \\
\midrule
\multicolumn{3}{l}{\textbf{Companion Computer}} \\
Computer & NVIDIA Jetson Orin Nano 8\,GB & Ubuntu 20.04, ROS 1 Noetic \\
\midrule
\multicolumn{3}{l}{\textbf{Motion Capture}} \\
Markers & Vicon markers ($\times$4) & 14\,mm diameter, asymmetric constellation \\
\bottomrule
\end{tabularx}
\end{table}

\subsection{Platform Variant Comparison}\label{sec:platform-comparison-detailed}

Table~\ref{tab:platform-comparison-detailed} provides a detailed comparison between the baseline platform developed in this work and the improved variant developed by Pablo Kirby.

\begin{table}[h]
\centering
\caption{Detailed comparison of baseline and improved X-wing platform variants.}
\label{tab:platform-comparison-detailed}
\small
\begin{tabularx}{\textwidth}{lXX}
\toprule
Parameter & Baseline Platform (This Work) & Improved Platform (Kirby, 2025) \\
\midrule
\textbf{Aerodynamics} & & \\
Airfoil & NACA~0015 (symmetric) & AG25 (cambered) \\
Dihedral/Anhedral & $\pm\SI{45}{\degree}$ (alternating) & $\pm\SI{20}{\degree}$ (alternating) \\
Wing Span & \SI{450}{\milli\meter} & \SI{500}{\milli\meter} \\
Wing Chord & \SI{250}{\milli\meter} & \SI{230}{\milli\meter} \\
Wing Thickness & \SI{37.5}{\milli\meter} (15\% chord) & \SI{57.5}{\milli\meter} (25\% chord) \\
Total Wing Area & \SI{0.318}{\meter\squared} & \SI{0.345}{\meter\squared} \\
Aspect Ratio (per wing) & 1.8 & 2.17 \\
\midrule
\textbf{Structural} & & \\
Wing Material & Bambu Lab PLA Aero & Bambu Lab PLA Aero \\
Frame Material & Standard PLA & Standard PLA \\
Motor Tilt & $\pm\SI{5}{\degree}$ & $\pm\SI{5}{\degree}$ \\
Total Mass & \SI{2.5}{\kilogram} & \SI{2.8}{\kilogram} \\
\midrule
\textbf{Propulsion} & Identical & Identical \\
Motors & T-MOTOR VELOX~V2808 (×4) & T-MOTOR VELOX~V2808 (×4) \\
Propellers & HQProp 7×3.5×3 (×4) & HQProp 7×3.5×3 (×4) \\
Battery & 6S 2P 2800\,mAh & 6S 2P 2800\,mAh \\
\midrule
\textbf{Avionics} & Identical & Identical \\
Flight Controller & Kakute~H7 v1.5 & Kakute~H7 v1.5 \\
Companion & Jetson Orin Nano 8\,GB & Jetson Orin Nano 8\,GB \\
\midrule
\textbf{Control} & Identical & Identical \\
Control Architecture & INDI with aerodynamic compensation & INDI with aerodynamic compensation \\
Control Frequency & 300\,Hz inner \& outer & 300\,Hz inner \& outer \\
\bottomrule
\end{tabularx}
\end{table}

\section{Thrust Characterization Details}\label{sec:thrust-map-detailed}

\subsection{Test Procedure}

Static thrust measurements were performed using the following procedure:

\begin{enumerate}
\item A single T-MOTOR VELOX V2808 motor with HQProp 7×3.5×3 propeller was mounted on a Rokubi Mini six-axis force--torque sensor using a custom 3D-printed fixture.
\item The motor was powered through the same flight controller (Kakute H7) and ESC configuration used in flight to ensure identical DShot timing and behavior.
\item Battery configuration: Two 6S 1400\,mAh Li-Po batteries connected in parallel (6S 2P), matching the flight configuration.
\item Motor commands were swept from 0\% to 100\% throttle in 1\% increments.
\item At each command level, the motor was allowed to stabilize for 1\,s, then force data were recorded for 2\,s and averaged.
\item The test was repeated with all four motors installed on the vehicle (three fixed, one on force sensor) to capture realistic voltage sag effects.
\item A total of 975 samples were recorded across the full throttle range.
\end{enumerate}

\subsection{Power Configuration Considerations}

During the bench tests, we observed that the command-to-thrust relationship depends on how many motors are simultaneously powered from the battery due to load-dependent voltage sag and internal resistance effects.

In particular, running only a single motor from a single 6S pack produced a slightly different mapping than powering all four motors from the flight-representative configuration with two 6S packs in parallel (6S 2P). The voltage sag under multi-motor load shifts the thrust curve downward by approximately 5--8\% at high throttle.

Unless stated otherwise, the mapping used for control in flight was calibrated under the latter 6S 2P, four-motor-powered condition, and all subsequent thrust conversions refer to this configuration.

\subsection{Curve Fitting}

The measured thrust vs. motor command data were fitted with a second-order polynomial:
\begin{equation}
T(u) = a_2 u^2 + a_1 u + a_0,
\end{equation}
where $u \in [0,1]$ is the normalized command input and $T$ is the thrust in Newtons.

For the RPM-based mapping used in the thesis:
\begin{equation}
T(\text{RPM}) = 5.15\times 10^{-2}\,\text{RPM}^2 - 0.09\,\text{RPM},
\end{equation}
with $R^2 = 0.997$.

This empirical mapping was integrated into the control pipeline to convert desired thrust values from the INDI controller into per-motor DShot commands. The calibrated relationship improved the accuracy of total thrust estimation and allowed for energy-based efficiency analysis during flight tests.

\section{Software Architecture Details}\label{sec:software-detailed}

\subsection{System Architecture}

The onboard software stack consists of ROS~1 Noetic running on the Jetson Orin Nano companion computer. The control architecture is divided into three main components:

\begin{enumerate}
\item \textbf{Ground Station} (Remote PC):
\begin{itemize}
\item ROS master node
\item Trajectory generation and mission planning
\item Real-time visualization (RViz)
\item Manual override interface
\item Data logging (rosbag)
\end{itemize}

\item \textbf{Companion Computer} (Jetson Orin Nano):
\begin{itemize}
\item Geometric outer loop controller (300\,Hz)
\item INDI inner loop controller (300\,Hz)
\item State estimation (Vicon feedback processing)
\item MAVLink communication interface
\item Local telemetry logging
\end{itemize}

\item \textbf{Flight Controller} (Kakute H7):
\begin{itemize}
\item Modified Betaflight firmware
\item IMU data acquisition (8\,kHz)
\item Motor RPM telemetry (DShot bidirectional)
\item Battery voltage monitoring
\item Motor command execution (DShot600, 8\,kHz update rate)
\end{itemize}
\end{enumerate}

\subsection{Communication Protocols}

\paragraph{Jetson $\leftrightarrow$ Ground Station:}
\begin{itemize}
\item Protocol: ROS 1 TCP/UDP
\item Medium: Wi-Fi 802.11ac (5\,GHz)
\item Bandwidth: $\approx$50\,Mbps sustained
\item Latency: 5--15\,ms (local network)
\end{itemize}

\paragraph{Jetson $\leftrightarrow$ Kakute H7:}
\begin{itemize}
\item Protocol: MAVLink 2
\item Medium: UART (serial)
\item Baud rate: 921\,600\,bit/s
\item Update rate: 300\,Hz (synchronized with control loop)
\item Latency: $<$3\,ms
\end{itemize}

\subsection{State Estimation Pipeline}

The state estimation pipeline used the Vicon motion capture system to provide real-time position and attitude feedback to the controller. In some experiments, an onboard Extended Kalman Filter (EKF) fused IMU data for state estimation backup.

The Vicon system and onboard IMU operated in a consistent right-handed coordinate convention:
\begin{itemize}
\item \textbf{X-axis}: Forward (red in figures)
\item \textbf{Y-axis}: Left
\item \textbf{Z-axis}: Up
\item \textbf{Euler angles}: Roll ($\phi$) about X, Pitch ($\theta$) about Y, Yaw ($\psi$) about Z
\end{itemize}

\subsection{Control Implementation}

All flight maneuvers, including point-to-point transitions and trajectory tracking, were executed fully autonomously from predefined setpoints. The control logic ran entirely on the Jetson companion computer, with the flight controller acting only as a sensor interface and motor command relay.

Manual override from the ground station was available for safety. Trajectories could be manually stopped and the vehicle landed or disarmed from the base station using a dedicated emergency stop button.

\section{Test Facility Details}\label{sec:facilities-detailed}

\subsection{Munich Indoor Arena}

\begin{itemize}
\item \textbf{Dimensions}: Approximately $7 \times 8$\,m flying volume
\item \textbf{Vicon system}: 8 cameras
\item \textbf{Sampling rate}: 200\,Hz
\item \textbf{Position accuracy}: $\pm 2$\,mm
\item \textbf{Safety features}: Safety nets, soft landing surface
\item \textbf{Use cases}: Agility experiments (3\,g circle maneuvers, aggressive transitions)
\end{itemize}

\subsection{AIDA Hall (Reutlingen University)}

\begin{itemize}
\item \textbf{Dimensions}: Approximately $60 \times 45$\,m flying volume
\item \textbf{Vicon system}: 24 cameras
\item \textbf{Sampling rate}: 200\,Hz
\item \textbf{Position accuracy}: $\pm 2$\,mm
\item \textbf{Environment}: Indoor, still-air conditions, no wind sources
\item \textbf{Use cases}: Steady forward-flight efficiency trials, high-speed tests (up to 20\,m/s)
\end{itemize}

\subsection{Test Protocol Details}

\paragraph{Pre-flight:}
\begin{enumerate}
\item Battery charge verification ($>$95\% capacity)
\item Vicon marker visibility check
\item Ground station connection test
\item Emergency stop procedure rehearsal
\item Trajectory parameter verification
\end{enumerate}

\paragraph{Flight sequence:}
\begin{enumerate}
\item Autonomous take-off to 1.5\,m altitude
\item Hover stabilization (10\,s)
\item Trajectory execution (step response, circle, or straight line)
\item Return to hover
\item Autonomous landing
\end{enumerate}

\paragraph{Post-flight:}
\begin{enumerate}
\item Data integrity check (rosbag verification)
\item Timestamp synchronization validation
\item Outlier detection (position jumps, IMU spikes)
\item Visual trajectory inspection
\item Battery voltage log review
\end{enumerate}

\paragraph{Acceptance criteria:}
\begin{itemize}
\item No Vicon position jumps $>$50\,mm
\item No IMU saturation events
\item Timestamp monotonicity preserved
\item No external disturbances observed (visual confirmation)
\item Consistent trajectory shape across repeated runs
\end{itemize}

Each experiment was repeated 3--5 times until these criteria were met for at least 3 consecutive runs. Only the best-quality datasets were used for analysis.

\section{Data Processing Pipeline}\label{sec:data-processing}

All experiment data were logged in rosbag format and post-processed using custom Python scripts. The processing pipeline consisted of:

\begin{enumerate}
\item \textbf{Data extraction}: rosbag $\rightarrow$ CSV conversion using \texttt{rosbag\_pandas}
\item \textbf{Timestamp alignment}: NTP-synchronized timestamps merged across Vicon and onboard telemetry
\item \textbf{Filtering}: Low-pass Butterworth filter (40\,Hz cutoff) applied to Vicon velocity estimates
\item \textbf{Thrust estimation}: Motor RPM $\rightarrow$ thrust conversion using calibrated mapping
\item \textbf{Coordinate transformations}: Body frame $\leftrightarrow$ world frame conversions using measured attitude
\item \textbf{Error calculation}: Trajectory tracking errors computed as $\|\mathbf{p}_d - \mathbf{p}\|$
\item \textbf{Statistical analysis}: Mean, RMS, max error, standard deviation computed for each run
\item \textbf{Visualization}: Matplotlib/pgfplots for publication-quality figures
\end{enumerate}

All processing scripts are documented in the \texttt{scripts/} directory of the thesis repository.

\section{Manufacturing Details}\label{sec:manufacturing}

\subsection{3D Printing Parameters}

\paragraph{Wing skins (PLA Aero):}
\begin{itemize}
\item \textbf{Printer}: Bambu Lab X1-Carbon
\item \textbf{Material}: Bambu Lab PLA Aero (lightweight variant)
\item \textbf{Print mode}: Vase mode (single-wall spiral)
\item \textbf{Layer height}: 0.2\,mm
\item \textbf{Wall thickness}: 0.4\,mm (single perimeter)
\item \textbf{Infill}: 0\% (hollow structure)
\item \textbf{Flow ratio}: 0.6 
\item \textbf{Print speed}: 100\,mm/s
\item \textbf{Nozzle temperature}: 245°C
\item \textbf{Bed temperature}: 60°C
\item \textbf{Support}: None required
\item \textbf{Print time}: $\approx$2.5\,h per wing
\end{itemize}

\paragraph{Frame components (standard PLA):}
\begin{itemize}
\item \textbf{Printer}: Bambu Lab X1-Carbon
\item \textbf{Material}: Bambu Lab PLA Basic (black)
\item \textbf{Layer height}: 0.2\,mm
\item \textbf{Wall lines}: 4 (1.6\,mm total)
\item \textbf{Infill}: 20\% grid
\item \textbf{Print speed}: 150\,mm/s
\item \textbf{Nozzle temperature}: 220°C
\item \textbf{Bed temperature}: 60°C
\item \textbf{Support}: Tree supports (auto-generated)
\end{itemize}